\documentclass [prd,reprint,eqsecnum,nofootinbib]  {revtex4-1}
\usepackage{amsfonts,amsmath,amscd,mathrsfs,xcolor,hyperref}
\newcommand{\lang}[1]{\textcolor{blue}{#1}}
\newcommand{\science}[1]{\textcolor{red}{#1}}
\newcommand{\remove}[1]{\textcolor{green}{#1}}
\def\SO{\mathsf{SO}}
\def\so{\mathsf{so}}
\def\SU{\mathsf{SU}}
\def\cP{{\cal P}}
\def\cD{{\cal D}}
\def\cC{{\cal C}}
\def\cL{{\cal L}}
\def\w{\omega}
\def\v{\wedge}
%\usepackage[colorinlistoftodos, shadow]{todonotes} %notatki: paczka
%Komendy do stosowania:
%\todo{Comment on margin with line}
%\todo[inline]{Comment in text}
%\todo[noline]{Comment on margin without line}
%\missingfigure{Missing picture here.}
%\listoftodos %notatki: lista rzeczy do poprawy

\begin{document}
\title{%
Local Maxwell symmetry and gravity}
\author{\science{Elian Vey}}
\affiliation{\science{Institute for Theoretical Cosmology, University of
Asterion, 1 Nebula Way, Asterion}}
\author{\science{Mira Solen}}
\affiliation{\science{Institute for Theoretical Cosmology, University of
Asterion, 1 Nebula Way, Asterion}}
\date{\today}
\begin{abstract}\noindent
\lang{We discuss physical aspects of a theory obtained by gauging the
AdS--Maxwell symmetry. The resulting theory combines Einstein gravity
with an $\sf{SO(3,1)}$ Yang--Mills field. A second tetrad can be
introduced and associated with a linear combination of the Lorentz
and Maxwell connections. Taking this tetrad as a fundamental
variable casts the theory into a first-order form of $f$--$g$
gravity.}
\science{We also analyze a homogeneous cosmological model, derive an
exact gauge-field solution, and present a numerical example.}
\end{abstract}
\maketitle

\section{Introduction}

The Maxwell symmetry is an extension of Poincar\'e symmetry and is a
\lang{symmetry of fields in a constant electromagnetic background}
\cite{Bacry:1970ye}, \cite{Schrader:1972zd}. This symmetry is of
interest \lang{from a purely algebraic perspective because it
circumvents a well-known theorem that forbids a central extension of}
this algebra (see e.g., \cite{Galindo:1967}, \cite{Soroka:2011tc},
\cite{Gomis:2009dm}).

\lang{Until recently, however, the Maxwell symmetry attracted
relatively little attention. This is somewhat surprising because
physical systems in constant electromagnetic backgrounds, or in
closely related settings, occur frequently in physics. This suggests
that the Maxwell symmetry may be physically relevant. For example,
canonical noncommutativity (for a review, see
\cite{Douglas:2001ba}) is ultimately related to quantum systems in
constant electromagnetic backgrounds and their generalizations.
Consequently, this symmetry may be relevant to noncommutative field
theories, as noted in \cite{Doplicher:1994tu}. It may also be
relevant to the quantum Hall effect.}

\lang{Setting these speculations aside, we investigate the
construction and properties of theories with local Maxwell
invariance. In Sec.~II we review the formulation in
\cite{Durka:2011nf} and derive its field equations. Sec.~III
discusses possible relations between gauged Maxwell theories and
bimetric gravity. Sec.~IV studies homogeneous cosmological solutions,
and Sec.~V summarizes the results.}

\section{Gauged Maxwell theory}

\lang{We begin with definitions and then review results from
\cite{Durka:2011nf}.}

The starting point is a generalization of Maxwell symmetry in which
the Poincar\'e subalgebra of the original Maxwell algebra
\cite{Bacry:1970ye}, \cite{Schrader:1972zd} is replaced by its
anti-de Sitter counterpart. \lang{Equivalently, we consider the
symmetry algebra of complex scalar and Dirac fields on anti-de Sitter
space in a constant electromagnetic background
\footnote{A detailed discussion will be presented in a forthcoming
paper.}. This algebra, isomorphic to
$\sf{so}(3,1)\oplus\sf{so}(3,2)$, was discussed previously in
\cite{Soroka:2006aj}, \cite{Bonanos:2010fw},
\cite{Lukierski:2010dy} and has the following form}
\cite{Durka:2011nf}
\begin{align}
&[\mathcal{P}_{a},\mathcal{P}_{b}]=i(\mathcal{M}_{ab}-\mathcal{Z}_{ab})\,
,\nonumber\\
&[\mathcal{M}_{ab},\mathcal{M}_{cd}]=-i(\eta_{ac}\mathcal{M}_{bd}+\ldots)\, ,\nonumber\\
&[\mathcal{M}_{ab},\mathcal{Z}_{cd}]=-i(\eta_{ac}\mathcal{Z}_{bd}+\ldots)\,,\nonumber\\
&[\mathcal{Z}_{ab},\mathcal{Z}_{cd}]=-i(\eta_{ac}\mathcal{Z}_{bd}+\ldots)\, ,\nonumber\\
&[\mathcal{M}_{ab},\mathcal{P}_c]=-i(\eta_{ac}\mathcal{P}_{b}-\eta_{bc}\mathcal{P}_{a}),\quad
[\mathcal{Z}_{ab},\mathcal{P}_{c}]=0\,,\label{1}
\end{align}
\lang{where $\ldots$ denotes three additional terms obtained by
antisymmetrizing the pairs of indices $(ab)$ and $(cd)$. The
original Maxwell algebra \cite{Soroka:2006aj,Bonanos:2010fw,
Lukierski:2010dy} follows from Eq.~(\ref{1}) by a contraction that
removes $\mathcal{M}_{ab}$ from the right-hand side of the first
commutator and sets the entire right-hand side of the fourth
commutator to zero.}

\lang{We gauge this algebra by introducing the connection}
\begin{equation}\label{2}
\mathbb{A}_{\mu}=\frac{1}{2}\omega_\mu{}^{ab}
\mathcal{M}_{ab}+\frac{1}{\ell}e_\mu{}^{a}
\mathcal{P}_{a}+\frac{1}{2}h_\mu{}^{ab} \mathcal{Z}_{ab}\,,
\end{equation}
and its curvature
\begin{equation}\label{3}
\mathbb{F}_{\mu\nu}=\partial_\mu\mathbb{A}_{\nu}-\partial_\nu\mathbb{A}_{\mu}-i[\mathbb{A}_{\mu},\mathbb{A}_{\nu}]\,,
\end{equation}
\lang{which decomposes into Lorentz, translational, and Maxwell
components:}
\begin{equation}\label{4}
\mathbb{F}_{\mu\nu}=\frac12\, F_{\mu\nu}^{ab}\, \mathcal{M}_{ab}
+\frac{1}{\ell}T_{\mu\nu}^{a}\,\mathcal{P}_{a}+ \frac12\,
G_{\mu\nu}^{ab}\, \mathcal{Z}_{ab}\,,
\end{equation}
where
\begin{eqnarray}
F_{\mu\nu}^{ab}&=&R^{ab}_{\mu\nu}+\frac{1}{\ell^2}( e^a_\mu e^b_\nu- e^a_\nu e^b_\mu) \,,\label{5a}\\
T_{\mu\nu}^{a}&=&D^\omega_\mu e^a_\nu -D^\omega_\nu e^a_\mu\,,\label{5b}\\
G_{\mu\nu}^{ab}&=&D^\omega_\mu h^{ab}_\nu -D^\omega_\nu
h^{ab}_\mu\nonumber-\frac{1}{\ell^2}( e^a_\mu e^b_\nu- e^a_\nu
e^b_\mu)\\ &+&(h^{ac}_\mu h^{\quad b}_{\nu\,c}-h^{ac}_\nu h^{\quad
b}_{\mu\,c})\, .\label{5c}
\end{eqnarray}
\lang{We define the Lorentz-covariant derivative by}
\begin{equation*}
D^\omega_\mu(\ast)\equiv
\partial_\mu(\ast)-i[\omega_\mu,(\ast)].
\end{equation*}

\lang{The gauge transformations of the fields $h_\mu{}^{ab}$,
$\omega_\mu{}^{ab}$, and $e_\mu{}^{a}$ are}
\begin{eqnarray}
        \delta_\Theta h^{ab}_{\mu}&=&D^\omega_\mu \tau^{ab}+
h^{ac}_\mu(\lambda+\tau)_{c}^{\;\,b}+h^{bc}_\mu(\lambda+\tau)^{a}_{\;\;c}\label{6}\quad\\
      \delta_\Theta \omega^{ab}_{\mu}&=& D^\omega_\mu \lambda^{ab}
      \label{7}\\
   \delta_\Theta e^{a}_{\mu}&=&-\lambda^a_{\;b}
    \,e^b_\mu\,,\label{8}
\end{eqnarray}
where $\lambda^{ab}$ and $\tau^{ab}$ denote the \lang{parameters of the local}
 Lorentz and Maxwell symmetries, respectively. The
\lang{curvature components transform covariantly:}
\begin{eqnarray}
        \delta_\Theta G^{ab}_{\mu\nu}&=&-[\tau,F_{\mu\nu}]^{ab}-[(\lambda+\tau),G_{\mu\nu}]^{ab}\label{9}\\
      \delta_\Theta F^{ab}_{\mu\nu}&=&-[\lambda,F_{\mu\nu}]^{ab}  \label{10}\\
    \delta_\Theta
    T^{a}_{\mu\nu}&=&-\lambda^a_{\;b}T^b_{\mu\nu}\,.\label{11}
\end{eqnarray}
\science{We have omitted the local translation transformations.
In first-order gravity, these are related to diffeomorphisms (up to
local Lorentz transformations and, in general, terms proportional to
the field equations).}

\science{A natural non-topological, Lorentz- and Maxwell-invariant
four-form built from the gauge fields is the Einstein--Hilbert action
with a cosmological term:}
\begin{equation}
S_{E} =\frac1{64\pi G}\,
\int\epsilon_{abcd}\left(R_{\mu\nu}{}^{ab}e_{\rho}^{c}e_{\sigma}^{d}
-\frac{\Lambda}{3}e_{\mu}^{a}e_{\nu}^{b}e_{\rho}^{ c}e_{\sigma}^{
d}\right)\epsilon^{\mu\nu\rho\sigma}\label{12}
\end{equation}
\lang{Additional curvature-squared invariants are possible; their
topological combinations do not affect the local equations of motion
(see \cite{Durka:2011nf} and references therein).}

We now comment on the relation between our construction and that in
\cite{deAzcarraga:2010sw}. The authors construct curvatures similar
to those in Eqs.~(\ref{5a})--(\ref{5c}), but use all possible
geometric combinations invariant under local Lorentz transformations
to build the action. Consequently, the theory in
\cite{deAzcarraga:2010sw} \lang{manifestly breaks Maxwell symmetry.
Its geometric action is therefore much richer than the one in}
\cite{Durka:2011nf} and presented above.

The Einstein action (\ref{12}) contains no Maxwell curvature $G$.
\science{Curvature-squared wedge-product invariants can be
topological, but metric contractions yield non-topological terms.
There are two independent Lorentz-invariant quadratic contractions
of the curvature of $\varpi$, corresponding to the two invariant
bilinear forms on $\mathsf{so}(3,1)$.}
\lang{The corresponding Yang--Mills action for the $\sf{SO(3,1)}$
gauge group was obtained in \cite{Soroka:2011tc,Durka:2011nf}:}
\begin{equation}\label{13}
 S_{1}= -  \frac1{4\varrho_1}\,\int d^4x\, e\, \big( F_{\mu\nu}{}^{ab} + G_{\mu\nu}{}^{ab}\big) \big( F^{\mu\nu}{}_{ab} +
    G^{\mu\nu}{}_{ab}\big)\, .
\end{equation}
\lang{Because $\epsilon^{abcd}$ is an invariant tensor of
$\sf{SO(3,1)}$, there is a second manifestly gauge-invariant action:}
\begin{equation}\label{13a}
 S_{2}= -  \frac1{4\varrho_2}\,\int d^4x\, e\, \big( F_{\mu\nu}{}^{ab} + G_{\mu\nu}{}^{ab}\big) \big( F^{\mu\nu}{}^{cd} +
    G^{\mu\nu}{}^{cd}\big) \epsilon_{abcd}\, .
\end{equation}
In the formulas above, $\varrho_1$ and $\varrho_2$ are
dimensionless coupling constants.

\lang{Using the explicit curvatures $F$ and $G$ in Eqs.~(\ref{5a})
and (\ref{5c}), these actions simplify as follows. In particular,}
$$
F_{\mu\nu}{}^{ab}(\omega, e) + G_{\mu\nu}{}^{ab}(h,\omega,e) =
H_{\mu\nu}{}^{ab}(\varpi)
$$
where $H_{\mu\nu}{}^{ab}(\varpi)$ is the curvature of
$\varpi_{\mu}{}^{ab}\equiv\omega_{\mu}{}^{ab}+h_{\mu}{}^{ab}$.
\lang{We henceforth treat $\varpi$ as an independent connection. The
full action is the sum of the Einstein action (\ref{12}) and the two
Yang--Mills terms:}
\begin{align}
S &= S_E + S_1 +S_2 \nonumber\\
&=\frac{1}{64\pi
G}\int\epsilon_{abcd}\left(R_{\mu\nu}{}^{ab}e_{\rho}^{c}e_{\sigma}^{d}
-\frac{\Lambda}{3}e_{\mu}^{a}e_{\nu}^{b}e_{\rho}^{ c}e_{\sigma}^{
d}\right)\epsilon^{\mu\nu\rho\sigma}\nonumber\\
&-\frac14\,\int  e\, \left(\frac1{\varrho_1}\, H_{\mu\nu}{}^{ab}
H^{\mu\nu}{}_{ab}+ \frac1{\varrho_2}\,H_{\mu\nu}{}^{ab}
H^{\mu\nu}{}^{cd}\,\epsilon_{abcd}\right)\, .\label{14}
\end{align}

\lang{The field equations resulting from this action include the
Einstein equations with a cosmological constant,}
\begin{equation}\label{15a}
    R_{\mu\nu}-\frac12\, g_{\mu\nu}\, R +\Lambda\, g_{\mu\nu}= 8\pi G T_{\mu\nu}
\end{equation}
whose right-hand side is the energy-momentum tensor of the field
$H$
\begin{align}
    T_{\mu\nu} &= \frac{1}{\varrho_1}\left(H_{\mu\lambda}{}^{ab} H_\nu{}^\lambda{}_{ab} -
    \frac14\, g_{\mu\nu}\,H_{\lambda\sigma}{}^{ab}
    H^{\lambda\sigma}{}_{ab}\right) \nonumber\\
    &+\frac{1}{\varrho_2}\left(H_{\mu\lambda}{}^{ab} H_\nu{}^\lambda{}^{cd} -
    \frac14\, g_{\mu\nu}\,H_{\lambda\sigma}{}^{ab}
    H^{\lambda\sigma}{}^{cd}\right)\epsilon_{abcd}\label{15}
\end{align}


The field equations for $\varpi$ differ slightly from the standard
Yang-Mills form and read
\begin{equation}\label{16}
    \nabla^\mu \left(\frac1{\varrho_1}\, H_{\mu\nu}{}^{ab} +\frac1{\varrho_2}\,\epsilon^{abcd}
    H_{\mu\nu}{}_{cd}\right)=0\,,
\end{equation}
Here, $\nabla$ acts on spacetime indices through the metric
connection determined by the tetrad $e^a_\mu$, and on algebra
indices through $\varpi$. \science{The Einstein--Cartan equation
obtained by varying the independent Lorentz connection sets the
spacetime torsion to zero in the absence of spin sources.}

Unfortunately, a gauge theory with a noncompact gauge group is
known to suffer from the unitarity problem at the quantum level
\cite{Hsu:1980zt}. Moreover, even classically, coupling such a theory to gravity is
problematic because its energy density is not
positive definite, as can be seen explicitly by inspecting the
$T_{00}$ component of the energy-momentum tensor in (\ref{15})
above. We discuss this issue in more detail in Sec.~IV, where we
investigate the properties of simple cosmological solutions of the
theory defined by the action (\ref{14}).

\section{Bi-metric interpretation}

\lang{We briefly discuss the relation between the gauged Maxwell
theory and bimetric theories. The gauge transformations of the
connection $\varpi$ introduced above have the same formal structure
as those of the gravitational connection $\omega$. We may therefore
introduce a second tetrad one-form $f^a$ and its torsion two-form
$\tau^a$:}
\begin{equation}\label{17}
    df^a + \varpi^a{}_b\wedge f^b=\tau^a\,.
\end{equation}
Having defined the new tetrad field $f^a$, we can construct another
gauge-invariant action, which is the Einstein action for the fields
$f$ and $\varpi$.
\lang{In particular,}
\begin{equation}
S_{E'} =\frac1{64\pi G'}\,
\int\epsilon_{abcd}\left(H_{\mu\nu}{}^{ab}f_{\rho}^{c}f_{\sigma}^{d}
-\frac{\Lambda'}{3}f_{\mu}^{a}f_{\nu}^{b}f_{\rho}^{ c}f_{\sigma}^{
d}\right)\epsilon^{\mu\nu\rho\sigma}\label{18}
\end{equation}
is manifestly gauge invariant, just as the Einstein action
(\ref{12}) is. The theory defined by the sum $S_E+S_{E'}$
\lang{strongly resembles the first-order $f$--$g$ theory of Isham,
Salam, and Strathdee} \cite{Isham:1971gm} (see
\cite{Banados:2008fi} for a recent discussion and references). A
contact term can be added between the two metrics
$g_{\mu\nu}=e_\mu{}^ae_\nu{}^b\eta_{ab}$ and
$q_{\mu\nu}=f_\mu{}^af_\nu{}^b\eta_{ab}$. \science{Unlike standard
bimetric models with two independent diffeomorphism symmetries, this
construction has a single spacetime diffeomorphism symmetry from the
outset. A metric contact term therefore does not reduce a
$(\mathrm{Diff})^2$ symmetry here; it is invariant under the
existing diagonal diffeomorphism symmetry.} The most general contact
term considered here is
\begin{equation}\label{19}
    S_{cont} =
   {M^2}(-g)^u(-q)^{(1/2-u)}\, V(g^{-1}q)
\end{equation}
\lang{where $(g^{-1}q)^\mu{}_\nu = g^{\mu\rho}q_{\rho\nu}$ and $V$ is
a scalar potential constructed from this tensor.}

\lang{Bimetric theories have been studied in a variety of contexts;
see, for example, \cite{Banados:2008fi,Damour:2002wu,Damour:2002ws,
ArkaniHamed:2002sp,Blas:2005yk,Speziale:2010cf} and the review
\cite{Clifton:2011jh}. Whether local Maxwell symmetry has a natural
role in brane-world models, Kaluza--Klein theories, or noncommutative
geometry is left for future work.}

\section{Example: a simple cosmological model}

\lang{We now study a homogeneous cosmological sector of the model
defined by the action (\ref{14}), including the issue of
non-positive energy discussed in Sec.~II.}

\lang{We use the spatially flat Friedmann--Robertson--Walker (FRW)
metric}
\begin{equation}\label{20}
    ds^2 =-N(t)^2dt^2+a(t)^2 (dx^2+dy^2+dz^2)
\end{equation}
\lang{We derive the reduced dynamics from the variational principle.
Substituting the metric (\ref{20}) into the Einstein action (\ref{12})
and suppressing the (formally infinite) comoving spatial volume gives}
\begin{equation}
S_E=\frac{1}{8\pi G}\int dt \left(-\frac{3}{N}a\dot a^2-\Lambda a^3
N\right)\,.\label{21}
\end{equation}

\lang{We next impose the spatial symmetries of the FRW background.
Rotationally invariant tensors are built from the Kronecker delta and
the three-dimensional Levi-Civita tensor. The nonzero components of
the connection $\varpi$ are therefore}
\begin{equation}\label{22}
    \varpi_i{}^{0\alpha} = f(t)\, \delta_i^\alpha\,,
    \quad \varpi_i{}^{\alpha\beta} = g(t)\,
    \epsilon_i{}^{\alpha\beta}\,,\quad \alpha,\beta=1,2,3\,.
\end{equation}
\lang{The appendix gives the symmetry-reduction argument for this ansatz.}

\lang{Substituting this ansatz and the metric (\ref{20}) into $S_1$
in Eq.~(\ref{13}) gives}
\begin{equation}
    S_1 =\frac1{\varrho_1}\int dt \left( \frac{3a}{N}\Big(\dot{g}^2 -\dot{f}^2\Big)
+ \frac{3N}{a} g^2\Big(4f^2-g^2\Big) \right)\label{23}
\end{equation}
\lang{The action $S_2$ in Eq.~(\ref{13a}) becomes}
\begin{equation}
    S_2 =\frac1{\varrho_2}\int dt \left( \frac{12a}{N}\, \dot f\,\dot g - \frac{24N}{a}\, f\, g^3\right)\label{24}
\end{equation}
\lang{For $\varrho_1>0$, the kinetic term for $f$ in Eq.~(\ref{23})
has the opposite sign from that for $g$, and the potential is
unbounded below along some field directions. These features reflect
the indefinite invariant form of the noncompact gauge algebra.}

\lang{Varying $S_E+S_1+S_2$ with respect to the scale factor $a(t)$
and lapse $N(t)$ gives the Friedmann equations. The lapse is varied
before choosing cosmic time, $N(t)=1$:}
\begin{align}
\left(\frac{\dot{a}}{a} \right)^2-\frac{\Lambda}{3}
 &=\frac{8\pi G}{3\varrho_1}\left[\frac{3}{a^2} \left(\dot g^2 - \dot
 f^2\right)
 -\frac{3}{a^4}\,g^2\left(4f^2-g^2\right)\right]\nonumber\\
 &+\frac{8\pi G}{3\varrho_2}\left[\frac{12}{a^2}\, \dot f \dot g +\frac{24}{a^4}\,f\, g^3\right]\,.\label{25}
\end{align}
\lang{The second equation is}
\begin{align}
\frac{2\ddot a}{a} &+ \left(\frac{\dot a}{a}\right)^2 - \Lambda
 =-\frac{8\pi G}{3\varrho_1}\left[\frac{3}{a^2} \left(\dot g^2 - \dot
 f^2\right)\right.\nonumber\\
 &-\left.\frac{3}{a^4}\,g^2\left(4f^2-g^2\right)\right]
 -\frac{8\pi G}{3\varrho_2}\left[\frac{12}{a^2}\, \dot f \dot g +\frac{24}{a^4}\,f\, g^3\right]\label{26}
\end{align}
\lang{Comparing these equations with the standard Friedmann
equations,}
$$
\left(\frac{\dot{a}}{a} \right)^2-\frac\Lambda3=\frac{8\pi
G}{3}\varepsilon\, , \quad \ddot a -\frac{\Lambda a}3=-\frac{4\pi
G}{3}(\varepsilon+3p)a
$$
\lang{we identify the energy density as}
\begin{align}
    \varepsilon&= \frac{1}{\varrho_1}\left[\frac{3}{a^2} \left(\dot g^2 - \dot
 f^2\right)
 -\frac{3}{a^4}\,g^2\left(4f^2-g^2\right)\right]\nonumber\\
 &+\frac{1}{\varrho_2}\left[\frac{12}{a^2}\, \dot f \dot g +\frac{24}{a^4}\,f\, g^3\right]\label{27}
\end{align}
\lang{which is not positive definite. The pressure is}
\begin{equation}\label{28}
    p=\frac\varepsilon3
\end{equation}
 \science{The relation $p=\varepsilon/3$ follows from the vanishing
trace of the four-dimensional Yang--Mills stress tensor. Consequently,
the gauge sector redshifts as radiation: $\varepsilon=\mathcal R/a^4$,
where $\mathcal R$ is a constant fixed by the gauge-field initial
data. Define $\mathcal K\equiv 8\pi G\mathcal R/3$. For
$\mathcal K>0$, the expanding branches are shown below; contracting
branches follow by time reversal:}
\begin{align*}
 &\remove{a(t)= C \sqrt{2 (t-t_0)}}\,, &\remove{\Lambda=0}\\
 &\remove{a(t)= C \sqrt{\cos \left(\sqrt{\frac{|\Lambda|}{3} } \left(2 (t-t_0)\right)\right)}}\,, &\remove{\Lambda<0}\\
 &\remove{a(t)= C \sqrt{\cosh \left(\sqrt{\frac{|\Lambda|}{3} } \left(2 (t-t_0)\right)\right)}}\,, &\remove{\Lambda>0}\\
 &\science{a^2(t)=2\sqrt{\mathcal K}\,(t-t_B)}\,, &\science{\Lambda=0}\\
 &\science{a^2(t)=\sqrt{\frac{3\mathcal K}{|\Lambda|}}\,
\sin\!\left(2\sqrt{\frac{|\Lambda|}{3}}\,(t-t_B)\right)}\,, &\science{\Lambda<0}\\
 &\science{a^2(t)=\sqrt{\frac{3\mathcal K}{\Lambda}}\,
\sinh\!\left(2\sqrt{\frac{\Lambda}{3}}\,(t-t_B)\right)}\,, &\science{\Lambda>0}
\end{align*}
\science{Here $t_B$ is the zero of $a^2$; for $\Lambda<0$ the sine
branch is restricted to intervals on which $a^2>0$ and describes a
recollapsing universe. The first Friedmann equation requires
$\varepsilon\geq-\Lambda/(8\pi G)$ wherever $H^2\geq0$. Thus
$\varepsilon\geq0$ for $\Lambda=0$, while a positive cosmological
constant permits a negative energy density down to this lower bound.
This is a constraint on allowed homogeneous initial data, not a
general positivity theorem for the gauge theory.}

\lang{For completeness, the equations of motion for $g$ and $f$ are}
\begin{equation}
\frac1{\varrho_1}\left[ -(a\dot g)\dot{} + \frac{4}a\, gf^2 -
\frac{2}a\, g^3\right] -\frac1{\varrho_2}\left[ (2a\dot f)\dot{} +
\frac{12}a\, fg^2 \right]=0\label{29}
\end{equation}
and
\begin{equation}
\frac1{\varrho_1}\left[ (a\dot f)\dot{} + \frac{4}a\, g^2f \right]
-\frac1{\varrho_2}\left[ (2a\dot g)\dot{} + \frac{4}a\, g^3
\right]=0\label{30}
\end{equation}
\lang{On the gauge-field equations, stress-energy conservation
reduces to the continuity equation}
\begin{equation}\label{31}
\dot \varepsilon +3H(\varepsilon+p)= \dot \varepsilon
+4H\varepsilon=0\,,
\end{equation}
\lang{where we used Eq.~(\ref{28}). Substituting Eq.~(\ref{27}) into
Eq.~(\ref{31}) gives}
\begin{equation}\label{32}
    \dot g\, \mbox{\framebox{(\ref{29})}} +\dot f\, \mbox{\framebox{(\ref{30})}}=0\,.
\end{equation}



\lang{For a time-independent solution with $g\ne0$, Eq.~(\ref{30}) gives}
\[
f=\frac{\varrho_1}{\varrho_2}\,g.
\]
\lang{Eq.~(\ref{29}) then requires a nonzero constant solution to satisfy}
\[
\left(\frac{\varrho_1}{\varrho_2}\right)^2=-\frac1{4},
\]
which is impossible for real nonzero couplings. \remove{The assertion
that the same conclusion follows for all profiles with
$\dot f,\dot g\propto a^{-1}$ is not established here and is omitted.}
\science{This argument excludes only constant solutions with
$g\ne0$; the separate $g=0$ branch must be treated independently.}

\section{Analytic and numerical solutions}

\science{The absence of a constant solution with $g\ne0$ does not
imply that the gauge-field equations have no nontrivial solutions. Introduce
conformal time by $d\eta=dt/a(t)$ and denote $d/d\eta$ by a prime.
For the consistent invariant branch $f(\eta)=0$, Eqs.~(\ref{29}) and
(\ref{30}) reduce to}
\begin{equation}
\science{g''+2g^3=0.}
\end{equation}
\science{Its first integral is $g'^2+g^4=g_0^4$ for initial data
$g(\eta_0)=g_0$, $g'(\eta_0)=0$. An exact solution is}
\begin{equation}
\science{f(\eta)=0,\qquad
g(\eta)=g_0\,\operatorname{cn}\!\left(
\sqrt{2}\,g_0(\eta-\eta_0),\,\frac{1}{\sqrt{2}}\right),}
\end{equation}
\science{where $\operatorname{cn}(u,k)$ is the Jacobi elliptic cosine
with modulus $k$. On this branch, Eq.~(\ref{27}) gives
$\mathcal R=a^4\varepsilon=3g_0^4/\varrho_1$, so $\varrho_1>0$
provides a positive-radiation solution whose scale factor is given
above. This explicitly supplies a nontrivial, time-dependent
solution of the problem left open by the constant-field ansatz.}
\science{For completeness, the $g=0$ branch also integrates directly:
Eqs.~(\ref{29}) and (\ref{30}) give $f''=0$, hence
$f(\eta)=f_0+c(\eta-\eta_0)$. Its radiation constant is
$\mathcal R=-3c^2/\varrho_1$, illustrating explicitly that negative
homogeneous gauge-field energy is possible when $\varrho_1>0$.}

\science{For a numerical example with both gauge amplitudes nonzero,
set $\varrho_1=\varrho_2=1$. The conformal-time equations are}
{\color{red}
\begin{align}
f''&=\frac85\,g f^2-\frac{28}{5}\,f g^2,\nonumber\\
g''&=\frac45\,g f^2-\frac45\,f g^2-2g^3.
\end{align}}
\begin{equation}
\science{\mathcal R=
\frac{3}{\varrho_1}\left[g'^2-f'^2-g^2(4f^2-g^2)\right]
+\frac{12}{\varrho_2}\left(f'g'+2fg^3\right).}
\end{equation}
\science{We integrated this autonomous system with classical
fourth-order Runge--Kutta, step size $\Delta\eta=5\times10^{-5}$,
and initial data $f(0)=0.2$, $g(0)=1$, $f'(0)=g'(0)=0$. The
conserved quantity $\mathcal R=a^4\varepsilon$ for these data is
$\mathcal R=7.32$. The table reports the numerical trajectory
(rounded to the shown precision):}
\begin{center}
\color{red}
\begin{tabular}{c|rrrr|r}
$\eta$ & $f$ & $g$ & $f'$ & $g'$ & $\mathcal R$\\ \hline
0.00 & 0.200000000 & 1.000000000 & 0.000000000 & 0.000000000 & 7.32\\
0.05 & 0.198682565 & 0.997343611 & -0.052595037 & -0.106111377 & 7.32\\
0.10 & 0.194760785 & 0.989417491 & -0.103975195 & -0.210507798 & 7.32\\
0.15 & 0.188324405 & 0.976348712 & -0.152998445 & -0.311556643 & 7.32\\
0.20 & 0.179517057 & 0.958342327 & -0.198661156 & -0.407782876 & 7.32
\end{tabular}
\end{center}
\science{The constancy of $\mathcal R$ is a useful check on both the
integration and the radiation equation of state. These data have
$\mathcal R>0$ and therefore source an expanding flat-FRW solution
for $\Lambda=0$ (or the corresponding branches for nonzero
$\Lambda$). At $\eta=0.20$, halving the step size changes each
reported gauge-field variable by less than $4\times10^{-16}$ in
double precision. If we also set $8\pi G=1$, then
$\mathcal K=\mathcal R/3=2.44$ and the $\Lambda=0$ scale factor is
$a^2(t)=2\sqrt{2.44}\,(t-t_B)$.}

\section{Conclusions}

\science{We reviewed the gauged AdS--Maxwell construction and its
interpretation as Einstein gravity coupled to a noncompact
$\mathsf{SO}(3,1)$ Yang--Mills sector, with a possible second tetrad
providing a bimetric interpretation. In the homogeneous, spatially
flat cosmological sector, the Yang--Mills stress tensor is traceless,
so its energy density scales as $a^{-4}$; its sign is not fixed by
the field equations. We corrected the scale-factor solutions implied
by the Friedmann equation, ruled out constant solutions with
$g\ne0$ for real couplings, and exhibited both an
exact elliptic-function solution and a reproducible numerical
trajectory with nonzero $f$ and $g$. The latter conserves the
radiation constant $\mathcal R$. A broader study of stability and of
the allowed initial-data space remains necessary, particularly
because the noncompact gauge group permits negative-energy
configurations.}






\appendix*
\section{Cosmological gauge connection}

\lang{In this appendix we derive the symmetry-reduced form of the
gauge connection $\varpi_\mu{}^{ab}$ on a spatially flat FRW
background. Its spatial isometries are translations and rotations,
which form the three-dimensional Euclidean group. We raise and lower
spatial indices with $\delta_{ij}$:}
\begin{equation}\label{a1}
    [R_i, R_j] = \epsilon_{ijk}\, R_k, \, [R_i, T_j] = \epsilon_{ijk}\, T_k, \, [T_i, T_j]
    =0\, .
\end{equation}
The Killing vector fields of the flat FRW spacetime form a
representation of this algebra
\begin{align}
T_i &= \frac{\partial}{\partial x^i}\nonumber \\ R_i
&=\epsilon_{ijk}\, x^j \,\frac{\partial}{\partial x^k}\label{a2}
\end{align}

The gauge-field invariance under spacetime symmetries means that the
Lie derivatives along the Killing vectors vanish up to gauge
transformations \cite{Forgacs:1979zs,Jackiw:1979ub}:
\begin{align}
    {\cal L}_{T_i}\,\varpi_\mu^{ab} &= \partial_\mu \lambda^{(T)}_i{}^{ab}
    + [\varpi_\mu,\lambda^{(T)}_i]^{ab}\,,\label{a3}\\  {\cal L}_{R_i}\,\varpi_\mu^{ab} &= \partial_\mu \lambda^{(R)}_i{}^{ab}
    + [\varpi_\mu,\lambda^{(R)}_i]^{ab}\,.\label{a4}
\end{align}
\lang{These equations satisfy integrability conditions that follow
from the identity}
$$
[{\cal L}_X, {\cal L}_Y] = {\cal L}_{[X,Y]}\, ,
$$
which gives the following consistency condition on the gauge parameters
$\lambda$
\begin{equation}\label{a5}
{\cal L}_X\lambda^{(Y)}-{\cal L}_Y\lambda^{(X)}
-[\lambda^{(X)},\lambda^{(Y)}] =\lambda^{([X,Y])}\,.
\end{equation}
For two translations, Eq.~(\ref{a5}) becomes
\begin{equation}\label{a6}
    \partial_i\lambda^{(T)}_j{}^{ab}-
    \partial_j\lambda^{(T)}_i{}^{ab} -
    [\lambda^{(T)}_i,\lambda^{(T)}_j]^{ab}=0\, .
\end{equation}
\science{This is the flatness condition for the compensating
translation gauge parameters. Locally, on a simply connected spatial
patch, a flat connection is pure gauge; choose that gauge so that
$\lambda^{(T)}_i{}^{ab}=0$. This is a gauge choice, not a conclusion
that the compensators vanish in every gauge.}

With this choice, Eq.~(\ref{a5}) for rotations and translations
reduces to the condition
$$
\partial_i\lambda^{(R)}_j{}^{ab}=0
$$
so $\lambda^{(R)}_j{}^{ab}$ is spatially constant. For two
rotations, Eq.~(\ref{a5}) becomes
\begin{equation}\label{a7}
    [\lambda^{(R)}_i,\lambda^{(R)}_j]^{ab}=\epsilon_{ijk}\,
    \lambda^{(R)}_k{}^{ab}\,,
\end{equation}
\lang{Thus the matrices $\lambda^{(R)}_i{}^{ab}$ form a
representation of $\sf{SO(3)}$.}

We can now return to Eqs.~(\ref{a3}) and (\ref{a4}), which reduce to
\begin{align}
    {\cal L}_{T_i}\,\varpi_\mu^{ab} &= 0\label{a8}\\  {\cal L}_{R_i}\,\varpi_\mu^{ab} &=  [\varpi_\mu,\lambda^{(R)}_i]^{ab}\label{a9}
\end{align}
\lang{Equation (\ref{a8}) implies that the connection depends only
on time, $\varpi_\mu^{ab}(x,t)=\varpi_\mu^{ab}(t)$. Equation
(\ref{a9}) then requires $\varpi_0^{ab}$ to commute with every
rotation generator. The Lorentz algebra has no nonzero element that
commutes with all spatial rotations, so $\varpi_0^{ab}=0$. For the
spatial components,}
\begin{equation}\label{a10}
\epsilon_{ijk}\, \varpi_j^{ab}=[\varpi_k,\lambda^{(R)}_i]^{ab}
\end{equation}
\lang{We choose the standard embedding of the rotation generators
into the Lorentz algebra $\sf{so}(3,1)$,}
\begin{equation}\label{a11}
\lambda^{(R)}_i{}^{\alpha\beta} = \epsilon_i{}^{\alpha\beta}\,,
\quad\lambda^{(R)}_i{}^{\alpha0} =0\,, \quad \alpha,\beta=1,2,3
\end{equation}
\lang{Substituting this representation into Eq.~(\ref{a10}) gives
two independent rotation-covariant structures:}
\begin{equation}\label{a12}
\varpi_j^{0\beta} =f(t) \, \delta_j^\beta\,, \quad \varpi_j^{\alpha\beta} =
g(t)\, \epsilon_j{}^{\alpha\beta}
\end{equation}
which reproduce the ansatz in the main text.

\begin{acknowledgments}
\remove{[REMOVE: the original acknowledgments and funding statement
refer to real individuals and grants and should not be attributed to
the fictional authors.]}
\end{acknowledgments}




















\begin{thebibliography}{99.}

%\cite{Bacry:1970ye}
\bibitem{Bacry:1970ye}
  H.~Bacry, P.~Combe, J.~L.~Richard,
  ``Group-theoretical analysis of elementary particles in an external electromagnetic field. 1.
  the relativistic particle in a constant and uniform field,''
  Nuovo Cim.\  {\bf A67 } (1970)  267-299.

%\cite{Schrader:1972zd}
\bibitem{Schrader:1972zd}
  R.~Schrader,
  ``The maxwell group and the quantum theory of particles in classical homogeneous electromagnetic fields,''
  Fortsch.\ Phys.\  {\bf 20 } (1972)  701-734.

  %\cite{Galindo:1967}
\bibitem{Galindo:1967}
  A.~Galindo,
  ``Lie algebra extensions of the Poincar\'e Algebra,''
  J.\ Math.\ Phys.\  {\bf 8 } (1967)  768.

%\cite{Soroka:2011tc}
\bibitem{Soroka:2011tc}
  D.~V.~Soroka and V.~A.~Soroka,
  ``Gauge semi-simple extension of the Poincar\'e group,''
  arXiv:1101.1591 [hep-th].
  %%CITATION = ARXIV:1101.1591;%%

%\cite{Gomis:2009dm}
\bibitem{Gomis:2009dm}
  J.~Gomis, K.~Kamimura and J.~Lukierski,
  ``Deformations of Maxwell algebra and their Dynamical Realizations,''
  JHEP {\bf 0908}, 039 (2009)
  [arXiv:0906.4464 [hep-th]].
  %%CITATION = JHEPA,0908,039;%%

  %\cite{Douglas:2001ba}
\bibitem{Douglas:2001ba}
  M.~R.~Douglas, N.~A.~Nekrasov,
  ``Noncommutative field theory,''
  Rev.\ Mod.\ Phys.\  {\bf 73 } (2001)  977-1029.
  [hep-th/0106048].

  %\cite{Doplicher:1994tu}
\bibitem{Doplicher:1994tu}
  S.~Doplicher, K.~Fredenhagen, J.~E.~Roberts,
  ``The Quantum structure of space-time at the Planck scale and quantum fields,''
  Commun.\ Math.\ Phys.\  {\bf 172 } (1995)  187-220.
  [hep-th/0303037].

%\cite{Bonanos:2009wy}
\bibitem{Bonanos:2009wy}
  S.~Bonanos, J.~Gomis, K.~Kamimura, J.~Lukierski,
 ``Maxwell Superalgebra and Superparticle in Constant Gauge Backgrounds,''
  Phys.\ Rev.\ Lett.\  {\bf 104}, 090401 (2010).
  [arXiv:0911.5072 [hep-th]].

%\cite{deAzcarraga:2010sw}
\bibitem{deAzcarraga:2010sw}
  J.~A.~de Azcarraga, K.~Kamimura, J.~Lukierski,
  ``Generalized cosmological term from Maxwell symmetries,''
  Phys.\ Rev.\  {\bf D83}, 124036 (2011).
  [arXiv:1012.4402 [hep-th]].






  %\cite{Soroka:2006aj}
\bibitem{Soroka:2006aj}
  D.~V.~Soroka, V.~A.~Soroka,
  ``Semi-simple extension of the (super)Poincare algebra,''
  Adv.\ High Energy Phys.\  {\bf 2009 } (2009)  234147.
  [hep-th/0605251].

  %\cite{Bonanos:2010fw}
\bibitem{Bonanos:2010fw}
  S.~Bonanos, J.~Gomis, K.~Kamimura, J.~Lukierski,
  ``Deformations of Maxwell Superalgebras and Their Applications,''
  J.\ Math.\ Phys.\  {\bf 51 } (2010)  102301.
  [arXiv:1005.3714 [hep-th]].

  %\cite{Lukierski:2010dy}
\bibitem{Lukierski:2010dy}
  J.~Lukierski,
  ``Generalized Wigner-Inonu Contractions and Maxwell (Super)Algebras,''
  Proc.\ Steklov Inst.\ Math.\  {\bf 272} (2011) 1
  [arXiv:1007.3405 [hep-th]].
  %%CITATION = PSMMB,272,1;%%

  %\cite{Durka:2011nf}
\bibitem{Durka:2011nf}
  R.~Durka, J.~Kowalski-Glikman, M.~Szczachor,
  ``Gauged AdS-Maxwell algebra and gravity,''
[arXiv:1107.4728 [hep-th]].



  %\cite{Hsu:1980zt}
\bibitem{Hsu:1980zt}
  J.~P.~Hsu, M.~D.~Xin,
  ``Incompatibility Of Unitarity And Gauge Symmetry In The Sl(2,c) Yang-mills Field Theory,''
  Phys.\ Rev.\  {\bf D24 } (1981)  471.

  %\cite{Isham:1971gm}
\bibitem{Isham:1971gm}
  C.~J.~Isham, A.~Salam, J.~A.~Strathdee,
  ``F-dominance of gravity,''
  Phys.\ Rev.\  {\bf D3 } (1971)  867-873.

  %\cite{Banados:2008fi}
\bibitem{Banados:2008fi}
  M.~Banados, A.~Gomberoff, D.~C.~Rodrigues, C.~Skordis,
  ``A Note on bigravity and dark matter,''
  Phys.\ Rev.\  {\bf D79 } (2009)  063515.
  [arXiv:0811.1270 [gr-qc]].

  %\cite{Damour:2002wu}
\bibitem{Damour:2002wu}
  T.~Damour, I.~I.~Kogan, A.~Papazoglou,
  ``Nonlinear bigravity and cosmic acceleration,''
  Phys.\ Rev.\  {\bf D66 } (2002)  104025.
  [hep-th/0206044].

  %\cite{Damour:2002ws}
\bibitem{Damour:2002ws}
  T.~Damour, I.~I.~Kogan,
 ``Effective Lagrangians and universality classes of nonlinear bigravity,''
  Phys.\ Rev.\  {\bf D66 } (2002)  104024.
  [hep-th/0206042].


%\cite{ArkaniHamed:2002sp}
\bibitem{ArkaniHamed:2002sp}
  N.~Arkani-Hamed, H.~Georgi, M.~D.~Schwartz,
  ``Effective field theory for massive gravitons and gravity in theory space,''
  Annals Phys.\  {\bf 305 } (2003)  96-118.
  [hep-th/0210184].

  %\cite{Blas:2005yk}
\bibitem{Blas:2005yk}
  D.~Blas, C.~Deffayet, J.~Garriga,
 ``Global structure of bigravity solutions,''
  Class.\ Quant.\ Grav.\  {\bf 23 } (2006)  1697-1719.
  [hep-th/0508163].

%\cite{Speziale:2010cf}
\bibitem{Speziale:2010cf}
  S.~Speziale,
  ``Bi-metric theory of gravity from the non-chiral Plebanski action,''
  Phys.\ Rev.\  {\bf D82 } (2010)  064003.
  [arXiv:1003.4701 [hep-th]].

%\cite{Clifton:2011jh}
\bibitem{Clifton:2011jh}
  T.~Clifton, P.~G.~Ferreira, A.~Padilla, C.~Skordis,
 ``Modified Gravity and Cosmology,''

  [arXiv:1106.2476 [astro-ph.CO]].




















%\cite{Forgacs:1979zs}
\bibitem{Forgacs:1979zs}
  P.~Forgacs, N.~S.~Manton,
  ``Space-Time Symmetries in Gauge Theories,''
  Commun.\ Math.\ Phys.\  {\bf 72 } (1980)  15.

%\cite{Jackiw:1979ub}
\bibitem{Jackiw:1979ub}
  R.~Jackiw, N.~S.~Manton,
  ``Symmetries And Conservation Laws In Gauge Theories,''
  Annals Phys.\  {\bf 127 } (1980)  257.

\end{thebibliography}


\end{document}
